\chapter{État de l’art}
\label{chap:etat_art}

\section{Classification automatique des stades du sommeil}

Les premières approches utilisaient des caractéristiques extraites manuellement (spectre de puissance, entropie, etc.) et des classifieurs classiques (SVM, forêts aléatoires). Depuis 2016, les réseaux de neurones profonds dominent le domaine.

\subsection{Modèles séquentiels : LSTM, CNN, Transformers}
\begin{itemize}
\item Les \textbf{CNN} sont efficaces pour extraire des motifs temporels locaux (fuseaux, ondes lentes).
\item Les \textbf{LSTM} (ou BiLSTM) modélisent les dépendances à long terme entre épisodes successifs.
\item Les \textbf{Transformers} utilisent des mécanismes d’auto-attention pour capturer des relations globales.
\end{itemize}
Plusieurs travaux combinent CNN et LSTM (DeepSleepNet, TinySleepNet) pour atteindre une précision > 85\%. Les Transformers récents (e.g., SleepTransformer) obtiennent des performances supérieures à 90\% sur Sleep-EDF.

\subsection{Modèles d’ensemble}

Le \textit{stacking} consiste à combiner plusieurs modèles de base via un méta-apprenant. Il permet de réduire la variance et le biais. Des études ont montré que l’empilement de CNN, LSTM et Transformer améliore la robustesse.

\section{Explicabilité en IA (XAI)}

\subsection{SHAP (SHapley Additive exPlanations)}
Basé sur la théorie des jeux, SHAP attribue à chaque caractéristique une contribution à la prédiction. Il est applicable à tout modèle.

\subsection{Mécanismes d’attention}
Les poids d’attention dans les LSTM ou Transformers indiquent quels instants temporels sont les plus importants.

\subsection{Cartes de saillie}
Pour les CNN, les cartes de saillie (saliency maps) mettent en évidence les zones du signal les plus influentes.

\subsection{Application au sommeil}
Peu de travaux intègrent l’explicabilité dans la classification des stades du sommeil. Le projet XAI-Sleep se positionne comme une contribution originale en combinant stacking et explications multi-méthodes.